% answers_short.tex —— 2005 年·普通高等学校招生考试·重庆卷·文科 · 简短答案速查
% 由 tools/make_short_answers.py 自 answers.tex 的 \gksolhead/\gkanswer 提取，
% 并与 2005/2005_数学真题_重庆文科_答案要点.md 逐题比对后生成，勿手改（改源 MD 后重跑）。

\begin{gknotice}
\textbf{使用说明}：本节只给最终答案便于快速核对，完整推导见解析版。答案为编写者独立推导，\textbf{非官方标准答案}。
\end{gknotice}

\gkpart{选择题}
\begin{center}\begin{tabular}{ccccccccccc}
\toprule
题号 & 1 & 2 & 3 & 4 & 5 & 6 & 7 & 8 & 9 & 10 \\
\midrule
答案 & A & D & D & B & C & B & B & A & A & C \\
\bottomrule
\end{tabular}\end{center}

\gkpart{填空题}
\begin{itemize}[leftmargin=2.6em]
  \item[\textbf{11.}] $(2, 3)$（或 $\{x \mid 2 < x < 3\}$）
  \item[\textbf{12.}] $\dfrac{8}{3}$
  \item[\textbf{13.}] $1$
  \item[\textbf{14.}] $2\sqrt{2}$
  \item[\textbf{15.}] $\dfrac{17}{45}$
  \item[\textbf{16.}] $x^2 + \dfrac{4y^2}{3} = 1$（或 $\dfrac{x^2}{1} + \dfrac{y^2}{\frac{3}{4}} = 1$）
\end{itemize}

\gkpart{解答题}
\begin{itemize}[leftmargin=2.6em]
  \item[\textbf{17.}] $a = \pm\sqrt{3}$。
  \item[\textbf{18.}] (1) $\dfrac{7}{10}$（或 $0.7$）；(2) 恰好取到一件合格品的概率为 $\dfrac{189}{1000}$；至少取到一件合格品的概率为 $\dfrac{973}{1000}$
  \item[\textbf{19.}] (1) $a = 3$；(2) $[0, +\infty)$（或 $a \geqslant 0$）
  \item[\textbf{20.}] (1) $1$；(2) $\dfrac{\pi}{4}$（或 $45^\circ$）
  \item[\textbf{21.}] (1) $\dfrac{x^2}{3} - y^2 = 1$；(2) $\left(-1, -\dfrac{\sqrt{3}}{3}\right) \cup \left(\dfrac{\sqrt{3}}{3}, 1\right)$
  \item[\textbf{22.}] (1) $b_1 = 2$，$b_2 = \dfrac{8}{3}$，$b_3 = 4$，$b_4 = \dfrac{20}{3}$；(2) $b_n = \dfrac{2^n + 4}{3}$；$S_n = \dfrac{2^n + 5n - 1}{3}$
\end{itemize}
